% Part: normal-modal-logic
% Chapter: tableaux
% Section: introduction

\documentclass[../../../include/open-logic-section]{subfiles}

\begin{document}

\olfileid[tr]{nml}{tab}{int}

\olsection{Giriş}

!!^{tableau}s, belirli (aşağıya doğru dallanan) !!{signed formula}s
ağaçlarıdır; başka bir deyişle, bir doğruluk değeri işareti ($\True$ veya
$\False$) ile !!a{sentence}{}den oluşan
\[
\sFmla{\True}{!A} \text{ veya } \sFmla{\False}{!A}
\]
çiftlerinin ağaçlarıdır. !!^a{tableau}, birtakım \emph{varsayımlarla}
başlar. Sonraki her !!{signed formula}, çıkarım kurallarından biri uygulanarak
üretilir. Bazı çıkarım kuralları ağacın bir ucuna bir ya da daha fazla
!!{signed formula} ekler; bazılarıysa iki yeni uç ekleyerek iki dal oluşturur.
Kuralların ürettiği !!{signed formula}{}deki !!{formula}, kuralın uygulandığı
!!{signed formula}{}dekinden daha az karmaşıktır. Bir dal hem
$\sFmla{\True}{!A}$ hem $\sFmla{\False}{!A}$ içerdiğinde o dala
\emph{kapalı} deriz. !!a{tableau}{}daki her dal kapalıysa !!{tableau}{}nun tamamı
kapalıdır. Kapalı bir !!{tableau}, !!{tableau}{}yu başlatmak için kullanılan
!!{signed formula}s kümesinin sağlanamaz olduğunu gösteren !!a{derivation}
oluşturur. Bu, bir $\Proves$ bağıntısı tanımlamak için kullanılabilir:
$\Gamma \Proves !A$ ancak ve ancak $\Gamma_0 = \{!B_1,
\dots, !B_n\} \subseteq \Gamma$ olacak biçimde, şu varsayımlar için kapalı
!!{tableau} bulunan sonlu bir küme varsa:
\[
\{\sFmla{\False}{!A}, \sFmla{\True}{!B_1}, \dots, \sFmla{\True}{!B_n}\}.
\]

Modal mantıklar için hem !!{signed formula} kavramını genişletmeli hem de
\iftag{prvBox}{$\Box$\iftag{prvDiamond}{ ve
    $\Diamond$}}{$\Diamond$} işleçlerini kapsayan kurallar eklemeliyiz. Modal
!!{tableau}s içindeki !!{formula}s, bir işaretin ($\True$ veya $\False$) yanı
sıra $\sigma$ \emph{öneklerine} de sahiptir. Önekler pozitif tam sayıların
boş olmayan dizileridir; yani $\sigma \in (\PosInt)^* \setminus
\{\emptyseq\}$. Bu önekleri çevrelerindeki $\tuple{\ }$ olmadan yazar ve
tek tek !!{element}s{}ı $,$ yerine $.$ ile ayırırız. $\sigma$ bir önekse
$\sigma.n$, $\sigma \concat \tuple{n}$ demektir; örneğin
$\sigma = 1.2.1$ ise $\sigma.3$, $1.2.1.3$'tür. Dolayısıyla örneğin
\[
\sFmla{\True}{\Box !A \lif !A}[1.2]
\]
bir \emph{önekli !!{signed formula}}dır (kısaca \emph{önekli
  !!{formula}}).

Sezgisel olarak önek, !!a{tableau} dalındaki !!{formula}s{}i sağlayabilecek bir
modeldeki dünyayı adlandırır; $\sigma$ bir dünyayı adlandırıyorsa $\sigma.n$,
$\sigma$'nın adlandırdığı dünyadan erişilebilen bir dünyayı adlandırır.

\end{document}
